\begin{figure}[htbp]
	\centering
	\begin{tikzpicture}[scale=1.1]
		\def\ax{5.5}    % Q1 dimension (horizontal): increased for separability plane gap
		\def\ay{3.0}  % Q2 dimension (vertical): 2cm circle + 1cm gap = 3cm
		\def\sx{2.5}  % horizontal shift for back face (Q3=1)
		\def\sy{2.0}  % vertical shift for back face (Q3=1)
		
		% Green separability plane: parallelogram parallel to side faces (Q1=0 and Q1=1)
		% midway between left (Q1=0) and right (Q1=1) groups
		\fill[green!20, opacity=0.4]
			({\ax/2}, -0.5) --
			({\ax/2}, {\ay+0.5}) --
			({\ax/2+\sx}, {\ay+\sy+0.5}) --
			({\ax/2+\sx}, {\sy-0.5}) -- cycle;
		\draw[green!60!black, thick]
			({\ax/2}, -0.5) --
			({\ax/2}, {\ay+0.5}) --
			({\ax/2+\sx}, {\ay+\sy+0.5}) --
			({\ax/2+\sx}, {\sy-0.5}) -- cycle;
		\node[green!60!black, font=\scriptsize] at ({\ax/2+\sx/2}, {(\ay+\sy)/2}) {separability};
		
		% Front face circles (Q3=0): solid, outer + inner (proportional sizes)
		\foreach \x/\y/\label/\fill/\innerR in {0/0/000/blue!50/1.2, \ax/0/100/blue!50/0.8, 0/\ay/010/blue!50/1.0, \ax/\ay/110/blue!50/0.6} {
			\node (\label) [draw, circle, minimum size=2.0cm, inner sep=0pt, thick] at (\x,\y) {};
			\node[draw, circle, fill=\fill, minimum size=\innerR cm, thick] at (\x,\y) {};
			\node[below, font=\tiny, fill=white, rounded corners, inner sep=1pt] at (\x,\y) {\label};
		}
		
		% Back face circles (Q3=1): dashed, outer + inner (proportional sizes)
		\foreach \x/\y/\label/\fill/\innerR in {\sx/\sy/001/green!50/1.0, \sx+\ax/\sy/101/green!50/0.6, \sx/\sy+\ay/011/green!50/0.8, \sx+\ax/\sy+\ay/111/green!50/0.4} {
			\node (\label) [draw, circle, minimum size=2.0cm, inner sep=0pt, thick, dashed, opacity=0.6] at (\x,\y) {};
			\node[draw, circle, fill=\fill, minimum size=\innerR cm, thick, dashed, opacity=0.6] at (\x,\y) {};
			\node[below, font=\tiny, fill=white, opacity=0.8, rounded corners, inner sep=1pt] at (\x,\y) {\label};
		}
		
		% Connect adjacent basis states by outer circle circumferences
		% Front face (Q3=0) adjacencies
		\draw ($(000.center)!1.0cm!(100.center)$) -- ($(100.center)!1.0cm!(000.center)$);
		\draw ($(000.center)!1.0cm!(010.center)$) -- ($(010.center)!1.0cm!(000.center)$);
		\draw ($(100.center)!1.0cm!(110.center)$) -- ($(110.center)!1.0cm!(100.center)$);
		\draw ($(010.center)!1.0cm!(110.center)$) -- ($(110.center)!1.0cm!(010.center)$);
		% Back face (Q3=1) adjacencies
		\draw ($(001.center)!1.0cm!(101.center)$) -- ($(101.center)!1.0cm!(001.center)$);
		\draw ($(001.center)!1.0cm!(011.center)$) -- ($(011.center)!1.0cm!(001.center)$);
		\draw ($(101.center)!1.0cm!(111.center)$) -- ($(111.center)!1.0cm!(101.center)$);
		\draw ($(011.center)!1.0cm!(111.center)$) -- ($(111.center)!1.0cm!(011.center)$);
		% Between front and back faces (Q3 changes)
		\draw ($(000.center)!1.0cm!(001.center)$) -- ($(001.center)!1.0cm!(000.center)$);
		\draw ($(100.center)!1.0cm!(101.center)$) -- ($(101.center)!1.0cm!(100.center)$);
		\draw ($(010.center)!1.0cm!(011.center)$) -- ($(011.center)!1.0cm!(010.center)$);
		\draw ($(110.center)!1.0cm!(111.center)$) -- ($(111.center)!1.0cm!(110.center)$);
		
		% Dimensional axes (top-left origin, all pointing down or right)
		\draw[->, thick] (-3.0,\ay+1.0) -- (-0.5,\ay+1.0) node[right] {$Q_1$};
		\draw[->, thick] (-3.0,\ay+1.0) -- (-3.0,\ay-1.5) node[below] {$Q_2$};
		\draw[->, thick] (-3.0,\ay+1.0) -- (\ax/2 - 3.5,\sy+\ay+0.5) node[right] {$Q_3$};
		
	\end{tikzpicture}
	\caption{Three-qubit state in Dimensional Circular Notation. The cube structure is shown with upright faces: front face ($Q_3 = 0$) and back face ($Q_3 = 1$) positioned top-right, each circle labeled by its basis state $|Q_1Q_2Q_3\rangle$. Three-qubit DCN: Eight circles in a $2 \times 2 \times 2$ cube structure. Front face (solid): $Q_3 = 0$. Back face (dashed): $Q_3 = 1$.}
	\label{fig:dcn-3qubit}
\end{figure}
